\heading{数学物理方法（下）-17春-期中}
\noteinfo{题目来源于赛艇，答案由AI生成。}

\begin{question}[简单计算和回答]

设第一大题含 3 个小问：
\begin{enumerate}[label=(\arabic*),leftmargin=2.6em,itemsep=0.25em,topsep=0.2em]
\item 在弦的横振动问题中，如果弦受到一个与速度成正比的阻尼（比例系数为 $d$），以及和弦的位移成正比的回复力（比例系数为 $k$），写出此时弦的振动方程。
\item 试在直角坐标系、柱坐标系、球坐标系中对三维球对称线性谐振子的薛定谔方程分离变量：
\[
 i\hbar\frac{\partial \psi}{\partial t}=-\frac{\hbar^2}{2m}\nabla^2\psi+\frac12 m\omega^2(x^2+y^2+z^2)\psi.
\]
\item 求解本征值问题，并计算本征函数的模方：
\[
\begin{cases}
X''(x)+\lambda X(x)=0, & -l\le x\le l,\\
X'(-l)=0,\quad X'(l)=0.
\end{cases}
\]
\end{enumerate}

\begin{solution}

\begin{enumerate}[label=(\arabic*),leftmargin=2.6em,itemsep=0.25em,topsep=0.2em]
\item 若已把张力和线密度归一化为波速 $c$，则弦位移 $u(x,t)$ 满足
\[
 u_{tt}+d u_t-c^2u_{xx}+ku=0.
\]
若不归一化，也可写成 $\rho u_{tt}+d u_t-Tu_{xx}+ku=0$。

\item 分离变量写成 $\psi=\Phi e^{-iEt/\hbar}$。

直角坐标下再令 $\Phi=X(x)Y(y)Z(z)$，可得
\[
X''+\Bigl(\frac{2mE_x}{\hbar^2}-\frac{m^2\omega^2}{\hbar^2}x^2\Bigr)X=0,
\]
\[
Y''+\Bigl(\frac{2mE_y}{\hbar^2}-\frac{m^2\omega^2}{\hbar^2}y^2\Bigr)Y=0,
\]
\[
Z''+\Bigl(\frac{2mE_z}{\hbar^2}-\frac{m^2\omega^2}{\hbar^2}z^2\Bigr)Z=0,
\quad E_x+E_y+E_z=E.
\]

柱坐标下令 $\Phi=R(\rho)\Theta(\varphi)Z(z)$，得
\[
\Theta''+m^2\Theta=0,
\]
\[
Z''+\Bigl(\frac{2mE_z}{\hbar^2}-\frac{m^2\omega^2}{\hbar^2}z^2\Bigr)Z=0,
\]
\[
R''+\frac1\rho R'+\Bigl(\frac{2mE_\perp}{\hbar^2}-\frac{m^2}{\rho^2}-\frac{m^2\omega^2}{\hbar^2}\rho^2\Bigr)R=0,
\quad E_\perp+E_z=E.
\]

球坐标下令 $\Phi=R(r)Y_{\ell}^{m}(\theta,\varphi)$，则角向部分给出球谐函数，径向满足
\[
R''+\frac{2}{r}R'+\Bigl(\frac{2mE}{\hbar^2}-\frac{m^2\omega^2}{\hbar^2}r^2-\frac{\ell(\ell+1)}{r^2}\Bigr)R=0.
\]

\item 令 $\lambda=\mu^2\ge 0$。通解为 $X=A\cos \mu x+B\sin \mu x$。由边界条件
\[
X'(\pm l)=0
\]
可得两族本征对：
\[
\lambda_0=0,\qquad X_0(x)=1,
\]
\[
\lambda_n^{(c)}=\Bigl(\frac{n\pi}{l}\Bigr)^2,\qquad X_n^{(c)}(x)=\cos\frac{n\pi x}{l},\quad n=1,2,\dots,
\]
\[
\lambda_n^{(s)}=\Bigl(\frac{(n+\tfrac12)\pi}{l}\Bigr)^2,
\qquad X_n^{(s)}(x)=\sin\frac{(n+\tfrac12)\pi x}{l},\quad n=0,1,2,\dots.
\]
模方为
\[
\|X_0\|^2=\int_{-l}^{l}1\,dx=2l,
\qquad
\|X_n^{(c)}\|^2=\|X_n^{(s)}\|^2=l.
\]
故归一化后可取
\[
\frac{1}{\sqrt{2l}},\qquad \frac{1}{\sqrt l}\cos\frac{n\pi x}{l},
\qquad \frac{1}{\sqrt l}\sin\frac{(n+\tfrac12)\pi x}{l}.
\]
\end{enumerate}

\ansmath{\|X_0\|^2=\int_{-l}^{l}1\,dx=2l,
\qquad
\|X_n^{(c)}\|^2=\|X_n^{(s)}\|^2=l\quad \frac{1}{\sqrt{2l}},\qquad \frac{1}{\sqrt l}\cos\frac{n\pi x}{l},
\qquad \frac{1}{\sqrt l}\sin\frac{(n+\tfrac12)\pi x}{l}}

\end{solution}
\end{question}

\begin{question}[一维弦受迫振动的半无界定解问题]

求解
\[
\begin{cases}
\dfrac{\partial^2u}{\partial t^2}-a^2\dfrac{\partial^2u}{\partial x^2}=A\cos\omega t,
&0\le x<\infty,\ t\ge 0,\\
\left.\dfrac{\partial u}{\partial x}\right|_{x=0}=0,\\
u|_{t=0}=x^2,\qquad \left.\dfrac{\partial u}{\partial t}\right|_{t=0}=x.
\end{cases}
\]

\begin{solution}

采用偶延拓。令全轴上的初位移、初速度分别为
\[
f_e(x)=x^2,\qquad g_e(x)=|x|,
\]
则满足 Neumann 边界条件的解可由全轴达朗贝尔公式给出：
\[
 u(x,t)=\frac{f_e(x+at)+f_e(x-at)}{2}+\frac{1}{2a}\int_{x-at}^{x+at}g_e(\xi)\,d\xi
 +\int_0^t (t-\tau)A\cos\omega\tau\,d\tau.
\]
化简得
\[
 u(x,t)=x^2+a^2t^2+\frac{1}{2a}\int_{x-at}^{x+at}|\xi|\,d\xi+\frac{A}{\omega^2}(1-\cos\omega t).
\]
其中
\[
\int_{x-at}^{x+at}|\xi|\,d\xi
=\frac{(x+at)|x+at|-(x-at)|x-at|}{2}.
\]
因此也可以写成
\[
 u(x,t)=x^2+a^2t^2+
 \frac{(x+at)|x+at|-(x-at)|x-at|}{4a}
 +\frac{A}{\omega^2}(1-\cos\omega t),\qquad x>0,\ t>0.
\]

\ansmath{u(x,t)=x^2+a^2t^2+
 \frac{(x+at)|x+at|-(x-at)|x-at|}{4a}
 +\frac{A}{\omega^2}(1-\cos\omega t),\qquad x>0,\ t>0}

\end{solution}
\end{question}

\begin{question}[导热细杆]

一个长度为 $l$ 的均匀导热细杆，侧面散热忽略，$x=0$ 端单位时间单位面积从外界吸收热量 $q$，$x=l$ 端保持温度为 $0$，初始温度分布为 $x-\cos x$。求杆中的温度分布随时间的变化。

\begin{solution}

记热方程为
\[
 u_t=\kappa u_{xx},\qquad 0<x<l,
\]
边界条件写成
\[
 -K u_x(0,t)=q,\qquad u(l,t)=0,
\]
其中 $K$ 为导热系数。先取稳态项
\[
 v(x)=\frac{q}{K}(l-x),
\]
则 $v''=0$ 且满足边界条件。令 $w=u-v$，则
\[
 w_t=\kappa w_{xx},\qquad w_x(0,t)=0,\quad w(l,t)=0.
\]
这是混合边值问题，其本征函数为
\[
 X_n(x)=\cos\mu_n x,
\qquad \mu_n=\frac{(n+\tfrac12)\pi}{l},\quad n=0,1,2,\dots
\]
故解可写成
\[
 u(x,t)=\frac{q}{K}(l-x)+\sum_{n=0}^{\infty}c_n e^{-\kappa\mu_n^2 t}\cos(\mu_n x),
\]
其中系数由初值
\[
 u(x,0)=x-\cos x
\]
决定，即
\[
 c_n=\frac{2}{l}\int_0^l\Bigl[x-\cos x-\frac{q}{K}(l-x)\Bigr]\cos(\mu_n x)\,dx.
\]
这就是温度分布的标准展开式。

\ansmath{c_n=\frac{2}{l}\int_0^l\Bigl[x-\cos x-\frac{q}{K}(l-x)\Bigr]\cos(\mu_n x)\,dx}

\end{solution}
\end{question}

\begin{question}[空心导体圆柱的电势]

一个半径为 $a$ 的无穷长空心导体圆柱，分成互相绝缘的两半，一半电势为 $V$，另一半接地，求圆柱内的电势分布。

\begin{solution}

这是二维 Laplace 方程的圆域边值问题。若以极坐标 $(r,\theta)$ 表示，并假定边界取值为
\[
 u(a,\theta)=
 \begin{cases}
 V,&0<\theta<\pi,\\
 0,&\pi<\theta<2\pi,
 \end{cases}
\]
则其 Fourier 展开为
\[
 u(a,\theta)=\frac{V}{2}+\frac{2V}{\pi}\sum_{m=0}^{\infty}
 \frac{\sin( (2m+1)\theta)}{2m+1}.
\]
圆内调和延拓为
\[
 u(r,\theta)=\frac{V}{2}+\frac{2V}{\pi}\sum_{m=0}^{\infty}
 \frac{1}{2m+1}\Bigl(\frac{r}{a}\Bigr)^{2m+1}\sin((2m+1)\theta),
 \qquad 0\le r<a.
\]
若原图所示两半的分割方向不同，只需把上式中的正弦项换成相应的余弦项即可，本质完全相同。

\ansmath{u(r,\theta)=\frac{V}{2}+\frac{2V}{\pi}\sum_{m=0}^{\infty}
 \frac{1}{2m+1}\Bigl(\frac{r}{a}\Bigr)^{2m+1}\sin((2m+1)\theta),
 \qquad 0\le r<a}

\end{solution}
\end{question}

\clearpage

